\documentclass[11pt]{article}
\usepackage[T1]{fontenc}
\usepackage[utf8]{inputenc}
\usepackage{lmodern}
\usepackage{microtype}
\usepackage[a4paper,margin=1in]{geometry}
\usepackage{amsmath,amssymb,amsthm,mathtools}
\usepackage{booktabs,longtable,array}
\usepackage{enumitem}
\usepackage{xcolor}
\usepackage{hyperref}
\usepackage[nameinlink,noabbrev]{cleveref}
\usepackage{url}
\hypersetup{colorlinks=true,linkcolor=blue!45!black,citecolor=green!35!black,urlcolor=blue!55!black,pdftitle={A Special-Function Continuation of the Alaoglu--Erdos Program},pdfauthor={OpenAI GPT-5.6 Pro}}
\setlist{itemsep=2pt,topsep=4pt}
\emergencystretch=3em
\allowdisplaybreaks
\numberwithin{equation}{section}

\newtheorem{theorem}{Theorem}[section]
\newtheorem{proposition}[theorem]{Proposition}
\newtheorem{lemma}[theorem]{Lemma}
\newtheorem{corollary}[theorem]{Corollary}
\newtheorem{criterion}[theorem]{Criterion}
\newtheorem{conjecture}[theorem]{Conjecture}
\theoremstyle{definition}
\newtheorem{definition}[theorem]{Definition}
\newtheorem{problem}[theorem]{Problem}
\theoremstyle{remark}
\newtheorem{remark}[theorem]{Remark}
\newtheorem{warning}[theorem]{Warning}

\newcommand{\Q}{\mathbb{Q}}
\newcommand{\Z}{\mathbb{Z}}
\newcommand{\N}{\mathbb{N}}
\newcommand{\C}{\mathbb{C}}
\newcommand{\R}{\mathbb{R}}
\newcommand{\Qbar}{\overline{\mathbb{Q}}}
\newcommand{\Li}{\operatorname{Li}}
\newcommand{\ord}{\operatorname{ord}}
\newcommand{\Div}{\operatorname{Div}}
\newcommand{\dist}{\operatorname{dist}}
\newcommand{\e}{\mathrm{e}}
\newcommand{\dd}{\,\mathrm{d}}

\title{\textbf{A Special-Function Continuation of the Alaoglu--Erd\H{o}s Program}\\[2mm]
\large Transfer quotients, polylogarithmic Cartier spectra,\\
$q$-gamma arithmetic, and Barnes double-zeta resonance}
\author{Research continuation prepared for the attached unified report}
\date{21 August 2026}

\begin{document}
\sloppy
\maketitle

\begin{abstract}
The Alaoglu--Erd\H{o}s conjecture asks whether a real number $x$ for which $2^x$ and $3^x$ are integers must itself be an integer.  This report continues, rather than repeats, the attached unified report.  It develops a special-function program around five exact encodings: an entire gamma quotient that removes all known integral common zeros; Gauss and Hurwitz transfer quotients whose multipliers are $q^x$; the polylogarithm $\Li_{-x}$ as a simultaneous Cartier eigenfunction; a mixed $(2,3)$-Mahler series whose Mellin transform acquires a Barnes-type double-pole resonance; and a double-sine function with periods $\log 2$ and $\log 3$ whose two contiguous quotients become algebraic under the hypothesis.

The strongest unconditional structural result is an ``algebraic-base spectrum'' theorem obtained from the Six Exponentials Theorem.  Under a hypothetical nonintegral solution, a positive algebraic number $a$ has algebraic $a^x$ \emph{if and only if} $a$ lies in the divisible hull of the multiplicative group generated by $2$ and $3$.  In particular, $n^x$ is an integer for every $3$-smooth integer $n$, while $n^x$ is transcendental for every integer having a prime factor at least $5$.  This turns the search for a proof into a precise rank-amplification problem: a special-function identity need only force algebraicity of one further value, such as $5^x$.

No unconditional final contradiction is obtained.  The report gives complete proofs of the new reductions, corrects a tempting but invalid ``common Barnes divisor'' heuristic, proves the exact Mellin resonance expansion (including its two incommensurable oscillatory terms), derives unconditional irrational $q$-digamma consequences from known Pad\'e theory, and isolates four concrete lemmas whose proof would settle the conjecture.
\end{abstract}

\tableofcontents
\newpage

\section{Scope, non-duplication boundary, and status}
\label{sec:scope}

The inherited source was found at \texttt{Article/alaoglu-erdos-unified-report/alaoglu-erdos-unified-report.tex}.  Its section map contained 367 section-level headings.  The present continuation deliberately does not reproduce that report's gamma finite-difference development, Barnes--$G$ determinant asymptotics, continuation layers, or its pre-existing bibliography survey.

The purpose here is narrower and more experimental: to ask what the classical and modern special functions actually add once the elementary reductions and the usual Four Exponentials/Schanuel perspective are already known.  The guiding rule is to distinguish three levels throughout:
\begin{enumerate}[label=(\roman*)]
\item identities that merely repackage the original equality;
\item unconditional arithmetic consequences that genuinely use transcendence theory;
\item proposed rank-amplification statements which are presently unproved but would imply the conjecture.
\end{enumerate}

The requested conclusion is a proof.  The investigation below does not honestly reach one.  What it does reach is a sharper description of what any special-function proof must accomplish, together with several exact equivalences and a nontrivial unconditional classification theorem.  In particular, the report never treats numerical evidence, natural boundaries, irrational special values, or a double Mellin pole as though any one of them were already a contradiction.

\section{The elementary reduction and a decisive algebraic-base spectrum}
\label{sec:spectrum}

Assume
\begin{equation}
  M:=2^x\in\Z,\qquad N:=3^x\in\Z.
  \label{eq:MN}
\end{equation}
If $x<0$, then $0<2^x,3^x<1$, impossible for integers; $x=0$ is an allowed integral solution.  If $x=p/q>0$ is rational in lowest terms and $2^x=M$, then $2^p=M^q$, so unique factorization gives $q=1$.  Thus every nonintegral solution is irrational.  The Gelfond--Schneider theorem then shows that it is in fact transcendental: an algebraic irrational $x$ would make $2^x$ transcendental, contrary to \eqref{eq:MN}.

The following theorem is the main unconditional structural gain of the present continuation.  It applies more generally than the pair $(2,3)$.

\begin{definition}
For multiplicatively independent positive algebraic numbers $a,b$, let
\[
 \Div\langle a,b\rangle
 :=\{c\in\Qbar_{>0}: c^k=a^r b^s\text{ for some }k\ge1,\ r,s\in\Z\}
\]
be the divisible hull of the rank-two multiplicative group they generate.
\end{definition}

\begin{theorem}[Algebraic-base spectrum]
\label{thm:spectrum}
Let $x\in\R\setminus\Q$, and let $a,b\in\Qbar_{>0}$ be multiplicatively independent.  Suppose $a^x,b^x\in\Qbar$.  Then for every $c\in\Qbar_{>0}$,
\begin{equation}
 c^x\in\Qbar
 \quad\Longleftrightarrow\quad
 c\in\Div\langle a,b\rangle.
 \label{eq:spectrum}
\end{equation}
\end{theorem}

\begin{proof}
If $c^k=a^r b^s$, then $(c^x)^k=(a^x)^r(b^x)^s$ is algebraic, so the positive real $k$th root $c^x$ is algebraic.

Conversely, suppose that $c\notin\Div\langle a,b\rangle$.  Then $\log a,\log b,\log c$ are linearly independent over $\Q$: any rational relation, after clearing denominators and exponentiating, would put $c$ in the divisible hull.  Also $1,x$ are $\Q$-linearly independent.  The Six Exponentials Theorem applied to these two and three numbers says that at least one of
\[
 a,\ b,\ c,\ a^x,\ b^x,\ c^x
\]
is transcendental.  The first five are algebraic by assumption, hence $c^x$ is transcendental.  This proves the reverse implication.
\end{proof}

\begin{corollary}[Exact spectrum for a hypothetical counterexample]
\label{cor:integer-spectrum}
If $x$ is a nonintegral solution of \eqref{eq:MN}, then
\begin{align}
 c^x\in\Qbar
 &\Longleftrightarrow c\in\Div\langle2,3\rangle
 &&(c\in\Qbar_{>0}),\label{eq:alg-spectrum-23}\\
 r^x\in\Q
 &\Longleftrightarrow r=2^u3^v\text{ for some }u,v\in\Z
 &&(r\in\Q_{>0}),\label{eq:rat-spectrum-23}\\
 n^x\in\Qbar
 &\Longleftrightarrow n=2^a3^b\text{ for some }a,b\in\Z_{\ge0}
 &&(n\in\N).\label{eq:int-spectrum-23}
\end{align}
In the last line, the algebraic values are actually the integers $M^aN^b$; every other $n^x$ is transcendental.
\end{corollary}

This is considerably stronger than the statement that a third prime would cause trouble.  It classifies \emph{all} positive algebraic bases visible to the exponent $x$ inside $\Qbar$.

\begin{criterion}[Rank amplification]
\label{crit:rank-amplification}
Under \eqref{eq:MN}, the conjecture follows as soon as one can prove that $c^x$ is algebraic for a single positive algebraic $c\notin\Div\langle2,3\rangle$.  In particular, it suffices to force $5^x\in\Qbar$.
\end{criterion}

\begin{proof}
If $x$ were nonintegral, \Cref{thm:spectrum} would make $c^x$ transcendental.  Thus $x$ must be rational, and the elementary reduction makes it integral.
\end{proof}

A useful subsidiary fact is that
\begin{equation}
  \frac{\log2}{\log3}\quad\text{is transcendental}.
  \label{eq:logratio-trans}
\end{equation}
Indeed, it is irrational by unique factorization.  If it were algebraic, Gelfond--Schneider applied to $3^{\log2/\log3}=2$ would give a contradiction.  This fact will control the geometry of zero lattices and the convergence of oscillatory residue sums below.

\section{An entire gamma quotient that removes all known solutions}
\label{sec:entire-quotient}

The reciprocal gamma function is entire and vanishes simply at $0,-1,-2,\dots$.  Hence
\begin{equation}
 E_b(z):=\frac{1}{\Gamma(1-b^z)}
 \label{eq:Eb}
\end{equation}
is entire, and its zero set is
\begin{equation}
 Z(E_b)=\left\{\frac{\log m+2\pi i k}{\log b}:m\in\N,\ k\in\Z\right\}.
 \label{eq:Ezeros}
\end{equation}
The raw common-zero formulation $E_2(x)=E_3(x)=0$ includes every known solution $x\in\Z_{\ge0}$.  It is more useful to divide those zeros out exactly.

\begin{proposition}[Extraneous-zero quotient]
\label{prop:Qb}
For each integer $b\ge2$, the meromorphic expression
\begin{equation}
 Q_b(z):=\frac{\Gamma(-z)}{\Gamma(1-b^z)}
       =\frac{E_b(z)}{1/\Gamma(-z)}
 \label{eq:Qb}
\end{equation}
extends to an entire function.  It is nonzero at every $z\in\Z_{\ge0}$.  The Alaoglu--Erd\H{o}s conjecture is equivalent to
\begin{equation}
 Z(Q_2)\cap Z(Q_3)=\varnothing.
 \label{eq:Qcoprime}
\end{equation}
\end{proposition}

\begin{proof}
The only poles of $\Gamma(-z)$ occur at $z\in\Z_{\ge0}$.  At such a point $z=k$, the denominator $\Gamma(1-b^z)$ also has a simple pole because $1-b^k$ is a nonpositive integer; the two singularities cancel with a finite nonzero ratio.  Elsewhere the numerator is finite and the gamma function in the denominator has no zeros, so no pole remains.  This proves entireness and nonvanishing at the removed integral points.

If $Q_2(z)=Q_3(z)=0$, then $2^z$ and $3^z$ are positive integers, while $z$ is not a nonnegative integer.  Writing $z=u+iv$, the phases require $v\log2, v\log3\in2\pi\Z$.  By the irrationality of $\log2/\log3$, this forces $v=0$.  Thus a common zero is precisely a nonintegral real counterexample.  The converse is immediate.
\end{proof}

At a nonintegral real zero $x$ with $b^x=B\in\N$, one has
\begin{equation}
 Q_b'(x)=\Gamma(-x)(-1)^B B!\log b.
 \label{eq:Qderivative}
\end{equation}
Consequently a hypothetical common zero satisfies
\begin{equation}
 \frac{Q_2'(x)}{Q_3'(x)}
 =(-1)^{M-N}\frac{M!}{N!}\frac{\log2}{\log3}.
 \label{eq:derivative-ratio}
\end{equation}
This local invariant is transcendental by \eqref{eq:logratio-trans}.  It is not by itself contradictory: the same logarithmic ratio also occurs after normalization at genuine integral solutions of the unquotiented functions.

The ring of entire functions is a B\'ezout domain.  Thus \eqref{eq:Qcoprime} is also equivalent to the existence of entire $A,B$ with
\begin{equation}
 A(z)Q_2(z)+B(z)Q_3(z)=1.
 \label{eq:entire-bezout}
\end{equation}
This is an exact but nonconstructive restatement.  A productive analytic route would be a quantitative corona estimate for $|Q_2|+|Q_3|$ with sufficiently controlled growth.  Standard asymptotics of $1/\Gamma$ turn such an estimate into a lower bound for the simultaneous distances of $2^t$ and $3^t$ from the positive integers.  No available estimate is strong enough to rule out an exact zero.

\begin{warning}[Why the obvious Barnes-divisor picture is wrong]
The vertical zero progressions of $E_2$ and $E_3$ have imaginary periods $2\pi i/\log2$ and $2\pi i/\log3$.  They meet only on the real axis.  A hypothetical common zero therefore does \emph{not} generate a two-dimensional common zero lattice, and one cannot insert a Barnes double-gamma function as a common entire divisor.  The legitimate Barnes double object appears later, in the Mellin dual of a mixed Mahler series, not in the common divisor of \eqref{eq:Eb}.
\end{warning}

\section{Transfer quotients: a common special-function template}
\label{sec:transfer}

Many special functions have multiplication or distribution formulas of the shape
\begin{equation}
 \mathcal T_q F(s,\cdot)=q^{A(s)}F(s,\cdot),
 \label{eq:transfer-template}
\end{equation}
where $\mathcal T_q$ refines an argument into $q$ congruence classes.  Choosing $A(s)=s$, $-s$, or $1-s$ realizes the same exponential character
\begin{equation}
 \chi_x(q):=q^x.
 \label{eq:character}
\end{equation}
The hypothesis says that this character takes integer values at $q=2$ and $q=3$.  \Cref{thm:spectrum} says more: under a counterexample, its algebraic locus on positive algebraic inputs is exactly a rank-two divisible group.

This perspective separates two sorts of identities.

\begin{definition}
A transfer identity is \emph{rank preserving} if, starting with transfer parameters in a multiplicative subgroup $G$, every parameter produced by composition, inversion, and refinement remains in the divisible hull of $G$.  It is \emph{rank amplifying} if it forces algebraicity of a transfer multiplier at a parameter outside that hull.
\end{definition}

The standard multiplication formulas for gamma, Hurwitz zeta, Lerch transcendents, polylogarithms, and $q$-shifted factorials are rank preserving.  Their cocycle law is ultimately $\chi_x(mn)=\chi_x(m)\chi_x(n)$.  Starting only from $2$ and $3$, this law can produce $6^x$, $12^x$, $(2/3)^x$, roots of such values, and nothing at a genuinely independent base such as $5$.

This is not a theorem forbidding every special-function proof.  It is a precise warning against an important circularity: repeatedly applying multiplication formulas can create an impressive tower of algebraic-looking identities while never leaving $\Div\langle2,3\rangle$.  A successful proof must introduce an additive transformation, a monodromy constraint, an arithmetic value theorem, or another mechanism that genuinely raises multiplicative rank.

\section{Gauss multiplication as a gamma-theoretic detector}
\label{sec:gauss}

Gauss's multiplication formula gives, for $q\ge2$,
\begin{equation}
 \prod_{r=0}^{q-1}\Gamma\!\left(\frac{x+r}{q}\right)
 =(2\pi)^{(q-1)/2}q^{1/2-x}\Gamma(x).
 \label{eq:gauss-mult}
\end{equation}
Accordingly define the normalized gamma transfer quotient
\begin{equation}
 \mathfrak G_q(x):=
 (2\pi)^{(q-1)/2}q^{1/2}
 \frac{\Gamma(x)}{\displaystyle\prod_{r=0}^{q-1}\Gamma((x+r)/q)}.
 \label{eq:Gq}
\end{equation}
Then the special-function expression collapses exactly to
\begin{equation}
 \mathfrak G_q(x)=q^x.
 \label{eq:Gq-qx}
\end{equation}

\begin{theorem}[Gamma transfer spectrum]
\label{thm:gamma-spectrum}
Under a hypothetical nonintegral solution of \eqref{eq:MN}, and for every integer $q\ge2$,
\begin{equation}
 \mathfrak G_q(x)\in\Qbar
 \quad\Longleftrightarrow\quad
 q=2^a3^b\quad(a,b\in\Z_{\ge0}).
 \label{eq:gamma-spectrum}
\end{equation}
For $3$-smooth $q$ the value is the integer $M^aN^b$; for every other $q$ it is transcendental.
\end{theorem}

\begin{proof}
Combine \eqref{eq:Gq-qx} with \Cref{cor:integer-spectrum}.
\end{proof}

For $q=2$ and $3$, \eqref{eq:gauss-mult} becomes
\begin{align}
 \Gamma(x/2)\Gamma((x+1)/2)
 &=2^{1-x}\sqrt\pi\,\Gamma(x),\label{eq:dup-at-x}\\
 \Gamma(x/3)\Gamma((x+1)/3)\Gamma((x+2)/3)
 &=2\pi\,3^{1/2-x}\Gamma(x).
 \label{eq:trip-at-x}
\end{align}
Thus the hypothesis supplies algebraic normalized products on every refinement grid with denominator $2^a3^b$.  Refining first by $2$ then by $3$, or in the opposite order, gives the same denominator-$6$ product.  The resulting compatibility is exact but automatic; it records $(2^x)(3^x)=6^x$ and does not create a fifth base.

\begin{criterion}[Gamma rank-amplification criterion]
\label{crit:gamma5}
The Alaoglu--Erd\H{o}s conjecture would follow from any theorem which, under the algebraicity of $\mathfrak G_2(x)$ and $\mathfrak G_3(x)$ for real $x$, forces $\mathfrak G_q(x)$ to be algebraic for one integer $q$ having a prime factor at least $5$.
\end{criterion}

The smallest target is
\begin{equation}
 \mathfrak G_5(x)
 =(2\pi)^2\sqrt5\,
 \frac{\Gamma(x)}{\prod_{r=0}^4\Gamma((x+r)/5)}=5^x.
 \label{eq:G5}
\end{equation}
Under a counterexample this value is transcendental.  A future algebraic-independence theorem for gamma values at the three incommensurable refinement grids $2,3,5$ could therefore settle the problem.  Existing gamma-value conjectures are mostly organized around algebraic arguments or rational shifts; here the central argument $x$ is itself transcendental, which is a serious mismatch with the strongest available results.

\section{Hurwitz zeta and Lerch distribution relations}
\label{sec:hurwitz}

For $q\ge2$, Hurwitz zeta satisfies the distribution identity
\begin{equation}
 \sum_{r=0}^{q-1}\zeta\!\left(s,\frac{a+r}{q}\right)
 =q^s\zeta(s,a),
 \label{eq:hurwitz-distribution}
\end{equation}
initially for $\Re s>1$ and then by analytic continuation.  For $a=1$ and real nonintegral $x>0$, $\zeta(x)$ is finite and nonzero, so
\begin{equation}
 \mathfrak Z_q(x):=
 \frac{\sum_{r=0}^{q-1}\zeta(x,(1+r)/q)}{\zeta(x)}=q^x.
 \label{eq:Zq}
\end{equation}
The same algebraic-base spectrum therefore gives:

\begin{proposition}[Hurwitz transfer spectrum]
Under a hypothetical counterexample, $\mathfrak Z_q(x)$ is algebraic exactly for the $3$-smooth integers $q$, and is transcendental for every other integer $q\ge2$.
\end{proposition}

The shift identity
\begin{equation}
 \zeta(-x,a)-\zeta(-x,a+1)=a^x
 \label{eq:hurwitz-shift}
\end{equation}
gives a second interpretation.  The function of the argument $a\mapsto\zeta(-x,a)$ has algebraic unit increments precisely at algebraic points $a$ in $\Div\langle2,3\rangle$; at every other positive algebraic $a$, the increment is transcendental.  In particular, among positive integer arguments its increments are integers exactly at the $3$-smooth integers and transcendental elsewhere.

At integral $x=d\ge0$, one has
\begin{equation}
 \zeta(-d,a)=-\frac{B_{d+1}(a)}{d+1},
 \label{eq:hurwitz-bernoulli}
\end{equation}
so the entire dependence on $a$ collapses to a Bernoulli polynomial.  For nonintegral $x$, no such finite polynomial collapse occurs.  This suggests a special-function rigidity program:

\begin{problem}[Hurwitz finite-rank bridge]
Find a natural finite-dimensionality or determinant condition on the collection
\[
 \left\{\zeta\!\left(-x,\frac{r}{2^a3^b}\right):a,b\ge0,\ 1\le r\le2^a3^b\right\}
\]
which follows from $2^x,3^x\in\Z$ and which can occur only when $x\in\Z_{\ge0}$.
\end{problem}

The multiplication formulas alone do not supply such finite rank.  They give one averaging relation at each level, leaving exponentially many independent-looking values.  This is another manifestation of the rank-preserving barrier.

The Lerch transcendent
\[
 \Phi(z,s,a)=\sum_{n\ge0}\frac{z^n}{(n+a)^s}
\]
packages both the shift and congruence-class decompositions.  Its distribution relations again have multipliers $q^{-s}$ and therefore reproduce $q^x$ when $s=-x$.  The Lerch formulation is flexible enough to introduce an auxiliary algebraic variable $z$, but no known value theorem turns the two algebraic multipliers at $2$ and $3$ into one at a third independent base.

\section{The polylogarithm, Cartier operators, and holonomic collapse}
\label{sec:polylog}

Set
\begin{equation}
 F_x(z):=\Li_{-x}(z)=\sum_{n\ge1}n^x z^n,\qquad |z|<1.
 \label{eq:Fx}
\end{equation}
This single special function makes three aspects of the problem visible: its coefficients, its root-of-unity distribution laws, and the distinction between integral and nonintegral order.

\subsection{Cartier eigenvalues}
For a power series $f(z)=\sum_{n\ge0}a_nz^n$, define the Cartier extraction operator
\begin{equation}
 U_qf(z):=\sum_{n\ge0}a_{qn}z^n.
 \label{eq:Uq}
\end{equation}
Then
\begin{equation}
 U_qF_x=q^xF_x.
 \label{eq:Cartier-eigen}
\end{equation}
Thus \eqref{eq:MN} says that one and the same polylogarithm is a simultaneous eigenfunction of $U_2$ and $U_3$ with positive integer eigenvalues $M,N$.

The root-of-unity filter gives the equivalent distribution identity
\begin{equation}
 \sum_{j=0}^{q-1}\Li_{-x}(\zeta_q^jz)
 =q^{x+1}\Li_{-x}(z^q),
 \label{eq:polylog-distribution}
\end{equation}
where $\zeta_q=\e^{2\pi i/q}$.  For $0<z<1$ define
\begin{equation}
 \mathfrak L_q(x;z):=
 \frac{1}{q}\frac{\sum_{j=0}^{q-1}\Li_{-x}(\zeta_q^jz)}{\Li_{-x}(z^q)}=q^x.
 \label{eq:Lq}
\end{equation}
Again, under a counterexample this quotient is algebraic exactly at $3$-smooth $q$.

\subsection{D-finiteness is exactly integrality of the order}
For $d\in\Z_{\ge0}$,
\begin{equation}
 \Li_{-d}(z)=\frac{zA_d(z)}{(1-z)^{d+1}},
 \label{eq:eulerian-polylog}
\end{equation}
where $A_d$ is an Eulerian polynomial.  Hence integral $x$ makes $F_x$ rational.  The converse is stronger than mere nonrationality.

\begin{theorem}[Holonomic dichotomy]
\label{thm:holonomic}
For real $x>0$, the function $F_x(z)=\Li_{-x}(z)$ is D-finite over $\C(z)$ if and only if $x\in\Z_{\ge0}$; in that case it is rational.
\end{theorem}

\begin{proof}[Proof sketch with the monodromy mechanism]
The forward direction for integral $x$ is \eqref{eq:eulerian-polylog}.  For $x\notin\Z$, Jonqui\`ere's local expansion gives, with $z=\e^\mu$,
\begin{equation}
 \Li_{-x}(\e^\mu)
 =\Gamma(1+x)(-\mu)^{-1-x}
  +\sum_{k\ge0}\zeta(-x-k)\frac{\mu^k}{k!},
 \label{eq:jonquiere}
\end{equation}
for $|\mu|<2\pi$ on a chosen branch.  Continuing around $z=1$ creates the fractional-power germ, while subsequent turns around $z=0$ replace $\mu$ by $\mu+2\pi i m$.  The germs $(\mu+2\pi i m)^{-1-x}$, $m\in\Z$, span an infinite-dimensional vector space when $x\notin\Z$: a finite relation, expanded for large $\mu$, forces all power moments of finitely many distinct shifts to vanish, and a Vandermonde argument makes every coefficient zero.  A solution of a linear differential equation of finite order has a finite-dimensional monodromy orbit.  Hence $F_x$ cannot be D-finite.  This is the same obstruction underlying the non-holonomicity results of Flajolet--Gerhold--Salvy.
\end{proof}

Consequently, the original conjecture is equivalent to the following statement internal to the polylogarithm family:

\begin{criterion}[Cartier-to-holonomy bridge]
\label{crit:cartier-holonomy}
If $F_x$ has integer eigenvalues for both $U_2$ and $U_3$, then $F_x$ is D-finite.
\end{criterion}

The criterion is not presently known.  The eigenvalue equations \eqref{eq:Cartier-eigen} control only coefficients whose indices are multiplied by $2$ or $3$; they do not close the full $2$- or $3$-kernel.

\subsection{A precise regularity obstruction}
A sequence $(a_n)$ is $q$-regular if the vector space spanned by its $q$-kernel
\[
 \{(a_{q^en+r})_{n\ge0}:e\ge0,\ 0\le r<q^e\}
\]
is finite-dimensional.

\begin{proposition}
\label{prop:qregular}
For $x>0$, the sequence $(n^x)_{n\ge1}$ is $q$-regular over $\C$ for one (equivalently every) integer $q\ge2$ if and only if $x\in\Z_{\ge0}$.
\end{proposition}

\begin{proof}
For $x=d\in\Z_{\ge0}$, every subsequence $(q^en+r)^d$ is a polynomial in $n$ of degree at most $d$, so the kernel lies in the $(d+1)$-dimensional space spanned by $1,n,\ldots,n^d$.

If $x\notin\Z_{\ge0}$, the kernel contains, up to nonzero scalars, the sequences $(n+q^{-e})^x$ for all $e\ge1$.  Any finite linear relation among distinct shifts gives, after expanding at infinity,
\[
 \sum_j c_j\alpha_j^m=0\qquad(m=0,1,2,\ldots),
\]
because $\binom{x}{m}\ne0$ for every $m$.  The first finitely many equations form a Vandermonde system, forcing all $c_j=0$.  Thus the kernel has infinite rank.
\end{proof}

This proves that a route through automatic/regular sequences must establish a genuinely new finite-kernel statement; integer Cartier eigenvalues by themselves are not finite regularity.  Modern Cobham--Mahler theorems for multiplicatively independent bases become relevant only after such a closure has been obtained \cite{AdamczewskiBell,AdamczewskiFaverjonI,AdamczewskiFaverjonII}.

\section{A mixed $(2,3)$-Mahler series and exact Mellin resonance}
\label{sec:mahler-barnes}

The $3$-smooth part of \eqref{eq:Fx} is the lacunary series
\begin{equation}
 H_{M,N}(z):=\sum_{a,b\ge0}M^aN^b z^{2^a3^b}.
 \label{eq:HMN}
\end{equation}
When \eqref{eq:MN} holds, this is exactly
\[
 H_{M,N}(z)=\sum_{\substack{n\ge1\\p\mid n\Rightarrow p\in\{2,3\}}}n^xz^n,
\]
and it has nonnegative integer coefficients.

\begin{proposition}[Mixed Mahler equation]
\label{prop:mixed-mahler}
The function \eqref{eq:HMN} satisfies
\begin{equation}
 H_{M,N}(z)=z+M H_{M,N}(z^2)+N H_{M,N}(z^3)-MN H_{M,N}(z^6),
 \label{eq:mixed-mahler}
\end{equation}
or, with commuting dilation operators $T_qf(z)=f(z^q)$,
\begin{equation}
 (1-MT_2)(1-NT_3)H_{M,N}=z.
 \label{eq:factored-mahler}
\end{equation}
\end{proposition}

\begin{proof}
The terms with $a\ge1$ are $MH(z^2)$, those with $b\ge1$ are $NH(z^3)$, and those with both are counted twice and removed by $MNH(z^6)$.  The remaining term is $z$.
\end{proof}

\begin{proposition}[Natural boundary]
\label{prop:natural-boundary}
For positive integers $M,N$, the unit circle is a natural boundary for $H_{M,N}$.
\end{proposition}

\begin{proof}
The radius of convergence is $1$: along the support, $\log(M^aN^b)=O(\log(2^a3^b))$, so the $n$th-root growth tends to $1$.  The coefficients are integers.  The series is not rational, because its nonzero support has counting function $O((\log X)^2)$ and hence density zero, whereas the Skolem--Mahler--Lech theorem makes the zero set of an eventually linearly recurrent sequence a finite union of arithmetic progressions up to a finite set, so its nonzero support is eventually periodic; an infinite eventually periodic support has positive density.  The P\'olya--Carlson theorem now gives the natural boundary \\cite{Polya1921,Lech1953}.
\end{proof}

The natural boundary is rigorous but not discriminatory: it occurs for the genuine integral solutions as well as for every hypothetical pair $M,N$.  The Mellin transform is more informative.

Let
\[
 h_{M,N}(t):=H_{M,N}(\e^{-t}),\qquad
 \alpha:=\frac{\log M}{\log2},\quad
 \beta:=\frac{\log N}{\log3}.
\]
Termwise Mellin integration gives, for $\Re s>\max(\alpha,\beta)$,
\begin{equation}
 \int_0^\infty h_{M,N}(t)t^{s-1}\dd t
 =\frac{\Gamma(s)}{(1-M2^{-s})(1-N3^{-s})}.
 \label{eq:mellin-transform}
\end{equation}
The two pole lattices are
\begin{align}
 s&=\alpha+\frac{2\pi i k}{\log2},&k&\in\Z,\label{eq:poles2}\\
 s&=\beta +\frac{2\pi i \ell}{\log3},&\ell&\in\Z.\label{eq:poles3}
\end{align}
They meet if and only if $\alpha=\beta$, and then their only common point is the real point $s=x:=\alpha=\beta$.

\begin{theorem}[Double-pole resonance]
\label{thm:resonance}
For positive integers $M,N$, the following are equivalent:
\begin{enumerate}[label=(\roman*)]
\item $\log M/\log2=\log N/\log3$;
\item the Mellin transform \eqref{eq:mellin-transform} has a double pole;
\item it has a double pole on the real axis, at $s=x=\log M/\log2$.
\end{enumerate}
In the resonant case, for every $0<\eta<x$,
\begin{equation}
 h_{M,N}(t)=t^{-x}\left[A\log\frac1t+B+P_2\!\left(\log\frac1t\right)+P_3\!\left(\log\frac1t\right)\right]
 +O(t^{-x+\eta})
 \label{eq:resonance-asymptotic}
\end{equation}
as $t\downarrow0$, where
\begin{align}
 A&=\frac{\Gamma(x)}{\log2\,\log3},\label{eq:A}\qquad\\
 B&=\Gamma(x)\left[
 \frac{\psi(x)}{\log2\,\log3}
 +\frac{1}{2\log2}+\frac{1}{2\log3}
 \right],\label{eq:B}
\end{align}
and $P_2,P_3$ are real-analytic periodic functions of periods $\log2$ and $\log3$, respectively.  Explicitly,
\begin{align}
 P_2(u)&=\sum_{k\in\Z\setminus\{0\}}
 \frac{\Gamma(x+2\pi i k/\log2)}{\log2\,[1-\exp(-2\pi i k\log3/\log2)]}
 \exp\!\left(\frac{2\pi i k u}{\log2}\right),\label{eq:P2}\\
 P_3(u)&=\sum_{\ell\in\Z\setminus\{0\}}
 \frac{\Gamma(x+2\pi i \ell/\log3)}{\log3\,[1-\exp(-2\pi i \ell\log2/\log3)]}
 \exp\!\left(\frac{2\pi i \ell u}{\log3}\right).
 \label{eq:P3}
\end{align}
The series converge absolutely and locally uniformly.
\end{theorem}

\begin{proof}
The pole statement follows from \eqref{eq:poles2}--\eqref{eq:poles3}; a nonreal intersection would imply a rational ratio $\log2/\log3$.  In the resonant case put $s=x+v$.  The standard expansion
\[
 \frac1{1-\e^{-\omega v}}=\frac1{\omega v}+\frac12+O(v)
\]
shows that the Laurent coefficients of the double pole are $A/v^2+B/v$, with $A,B$ as stated.  Mellin inversion converts $A/v^2$ to $At^{-x}\log(1/t)$ and $B/v$ to $Bt^{-x}$.  The remaining poles on $\Re s=x$ give \eqref{eq:P2} and \eqref{eq:P3}.  Stirling's estimate gives exponential decay of $\Gamma(x+i\tau)$.  The possible small denominators are bounded below by a negative power of $|k|$ or $|\ell|$ using lower bounds for nonzero linear forms in $\log2$ and $\log3$; exponential gamma decay therefore gives absolute convergence.  Shifting the contour to $\Re s=x-\eta$ yields the stated remainder.
\end{proof}

The logarithmic term in \eqref{eq:resonance-asymptotic} is a clean analytic signature of the equality of the two exponents.  But it is not a proof of their integrality.  Integer-coefficient generating functions can have nonintegral singular exponents, and the two quasiperiodic corrections survive in both the integral and hypothetical nonintegral cases.

\subsection{Where Barnes double zeta really enters}
In the resonant case, put $\omega_1=\log2$, $\omega_2=\log3$, and $v=s-x$.  The geometric kernel is
\begin{equation}
 K(v)=\frac1{(1-\e^{-\omega_1v})(1-\e^{-\omega_2v})}
 =\sum_{a,b\ge0}\e^{-v(a\omega_1+b\omega_2)}.
 \label{eq:Barnes-kernel}
\end{equation}
This is precisely the heat kernel occurring in Barnes's double zeta function
\begin{equation}
 \zeta_2(w,a\mid\omega_1,\omega_2)
 =\frac1{\Gamma(w)}\int_0^\infty
 \frac{t^{w-1}\e^{-at}}{(1-\e^{-\omega_1t})(1-\e^{-\omega_2t})}\dd t.
 \label{eq:Barnes-zeta}
\end{equation}
The derivative at $w=0$ defines the Barnes double gamma function after normalization.  Thus a common exponent puts the Mellin resolvent into the canonical two-period Barnes class.  This is the correct two-dimensional structure: a lattice in the \emph{logarithms of the dilation semigroup}, not a common zero lattice of $E_2$ and $E_3$.

\begin{criterion}[Arithmetic resonance rigidity]
\label{crit:resonance-rigidity}
A theorem asserting that a mixed Mahler series of the form \eqref{eq:HMN} with positive integer multiplicities can have a real Barnes double-pole resonance only at an integer abscissa would prove the Alaoglu--Erd\H{o}s conjecture.
\end{criterion}

As stated, this criterion is exactly tailored to the problem.  Neither P\'olya--Carlson nor current multivariate Mahler theory supplies the required restriction on the pole location.

\section{A two-period double-sine coupling}
\label{sec:double-sine}

Choose the reciprocal-double-sine normalization
\[
 S_2(z\mid\omega_1,\omega_2)
 =\frac{\Gamma_2(z\mid\omega_1,\omega_2)}
 {\Gamma_2(\omega_1+\omega_2-z\mid\omega_1,\omega_2)},
\]
satisfies the two shift equations
\begin{align}
 \frac{S_2(z+\omega_1)}{S_2(z)}&=2\sin\frac{\pi z}{\omega_2},\label{eq:S2shift1}\\
 \frac{S_2(z+\omega_2)}{S_2(z)}&=2\sin\frac{\pi z}{\omega_1}.
 \label{eq:S2shift2}
\end{align}
Set
\begin{equation}
 \omega_1=\log2,\qquad \omega_2=\log3,\qquad
 z_x=\frac{i x\omega_1\omega_2}{2\pi}.
 \label{eq:zx}
\end{equation}
Then
\begin{align}
 2\sin\frac{\pi z_x}{\omega_2}
 &=i\left(2^{x/2}-2^{-x/2}\right)
 =i\left(M^{1/2}-M^{-1/2}\right),\label{eq:S2M}\\
 2\sin\frac{\pi z_x}{\omega_1}
 &=i\left(3^{x/2}-3^{-x/2}\right)
 =i\left(N^{1/2}-N^{-1/2}\right).
 \label{eq:S2N}
\end{align}
Thus the original hypothesis is encoded by the algebraicity of two adjacent shift quotients of a \emph{single} special function with incommensurable periods.

This is more genuinely two-base than separate gamma multiplication formulas.  Nevertheless, the commutativity of the two shifts creates no contradiction: traversing an elementary period parallelogram in either order gives the same result because the two sine factors acquire matching signs.  The compatibility is built into $S_2$.

\begin{problem}[Double-sine arithmetic rigidity]
Develop a value theorem for $S_2$ with logarithmic periods $\log2,\log3$ showing that simultaneous algebraicity of the two quotients in \eqref{eq:S2shift1}--\eqref{eq:S2shift2} at a point of the form \eqref{eq:zx} forces $x\in\Q$.
\end{problem}

Such a theorem would settle the conjecture.  Existing arithmetic results for multiple gamma and sine functions are not designed for a transcendental period ratio such as \eqref{eq:logratio-trans}.  The modular-double and hyperbolic-gamma literature does, however, provide analytic tools and transformation formulas that make this a concrete research direction rather than a purely formal reformulation \cite{Barnes1901,Ruijsenaars2000,Kurokawa1991,Narukawa2004}.

\section{$q$-gamma, $q$-digamma, and generalized $q$-logarithms}
\label{sec:qgamma}

For $0<q<1$, Jackson's $q$-gamma function is
\begin{equation}
 \Gamma_q(x)=(1-q)^{1-x}\frac{(q;q)_\infty}{(q^x;q)_\infty},
 \qquad (a;q)_\infty:=\prod_{k\ge0}(1-aq^k).
 \label{eq:qgamma}
\end{equation}
Take $q=1/b$ and suppose $B=b^x\in\N$.  Then $q^x=B^{-1}$, so
\begin{equation}
 \frac{(1-b^{-1})^{x-1}\Gamma_{1/b}(x)}{(b^{-1};b^{-1})_\infty}
 =\frac1{(B^{-1};b^{-1})_\infty}.
 \label{eq:qpoch-reduction}
\end{equation}
For $b=2,3$, the extra factor $(1-b^{-1})^x$ is itself rational under \eqref{eq:MN}: it is $M^{-1}$ for $b=2$ and $M/N$ for $b=3$.  Hence the hypothesis reduces two $q$-gamma values at the transcendental argument $x$ to $q$-Pochhammer products with rational parameters.

The logarithmic derivative is still more explicit:
\begin{equation}
 \psi_q(x)=-\log(1-q)+(\log q)\sum_{k\ge0}\frac{q^{k+x}}{1-q^{k+x}}.
 \label{eq:qdigamma}
\end{equation}
With $q=1/b$ and $B=b^x$,
\begin{equation}
 L_b(B):=\sum_{k\ge0}\frac1{Bb^k-1}
 =-\frac{\psi_{1/b}(x)+\log(1-b^{-1})}{\log b}.
 \label{eq:Lb-qdigamma}
\end{equation}
Expanding geometrically gives
\begin{equation}
 L_b(B)=\frac1{B-1}+\ell_b(B^{-1}),\qquad
 \ell_b(z):=\sum_{r\ge1}\frac{z^r}{b^r-1},
 \label{eq:qlog}
\end{equation}
and the generalized $q$-logarithm obeys the elementary $q$-difference equation
\begin{equation}
 \ell_b(z)-\ell_b(z/b)=\frac{z}{b-z}.
 \label{eq:qlog-difference}
\end{equation}

Borwein's Pad\'e theorem on series $\sum_{n\ge1}(q^n+r)^{-1}$ applies after the rescaling
\begin{equation}
 L_b(B)=\frac{b}{B}\sum_{n\ge1}\frac1{b^n-b/B}.
 \label{eq:Borwein-rescale}
\end{equation}
For integers $b,B\ge2$ these parameters are rational and nondegenerate, so $L_b(B)$ is irrational \cite{Borwein1992}.  We therefore obtain an unconditional consequence.

\begin{theorem}[$q$-digamma consequence]
\label{thm:qdigamma-irr}
If \eqref{eq:MN} holds with $x>0$, then both
\begin{equation}
 -\frac{\psi_{1/2}(x)+\log(1/2)}{\log2}=L_2(M),\qquad
 -\frac{\psi_{1/3}(x)+\log(2/3)}{\log3}=L_3(N)
 \label{eq:two-qdigamma}
\end{equation}
are irrational.
\end{theorem}

This is a genuine arithmetic theorem about special values at the unknown $x$.  It still does not distinguish an integral solution from a hypothetical nonintegral one: the same irrationality persists when $M=2^d,N=3^d$.  A proof would need a cross-base relation between the two $q$-logarithms or $q$-gamma products.  Standard $q$-multiplication formulas relate $q$ to its powers, not $1/2$ to $1/3$, and therefore remain rank preserving.

\section{Hypergeometric addition as a possible rank amplifier}
\label{sec:hypergeometric}

The simplest hypergeometric identity is
\begin{equation}
 {}_1F_0(-x;\,; -z)=(1+z)^x.
 \label{eq:1F0}
\end{equation}
It converts addition of algebraic bases into a special value.  Under a counterexample, \Cref{thm:spectrum} gives the complete arithmetic behavior at positive algebraic $z$:
\begin{equation}
 {}_1F_0(-x;\,;-z)\in\Qbar
 \quad\Longleftrightarrow\quad
 1+z\in\Div\langle2,3\rangle.
 \label{eq:hyper-spectrum}
\end{equation}
In particular,
\begin{align}
 {}_1F_0(-x;\,;-1)&=2^x=M,\label{eq:F0-2}\\
 {}_1F_0(-x;\,;-2)&=3^x=N,\label{eq:F0-3}\\
 {}_1F_0(-x;\,;-4)&=5^x\notin\Qbar.\label{eq:F0-5}
\end{align}
Equivalently,
\begin{equation}
 {}_1F_0(-x;\,;-2/3)=(5/3)^x=\frac{5^x}{N}
 \label{eq:additive-bridge}
\end{equation}
is forced to be transcendental.

This is the cleanest place where addition can raise multiplicative rank: $2+3=5$.  A proof strategy would seek a second evaluation of \eqref{eq:additive-bridge}, derived from the gamma/Hurwitz/$q$-gamma towers at bases $2$ and $3$, which makes it algebraic.  Classical quadratic and cubic hypergeometric transformations were examined for this purpose.  Their prefactors are themselves powers of algebraic rational functions; on closed transformation loops the prefactors multiply to $1$ identically in $x$.  Consequently, the standard transformations do not produce a hidden $5^x$ from $2^x$ and $3^x$.  Connection coefficients introduce gamma values at affine functions of $x$, returning the problem to \Cref{crit:gamma5} rather than resolving it.

A more sophisticated possibility is to use an algebraic transformation of a higher hypergeometric or Appell system whose monodromy representation has an arithmetic rigidity theorem.  The parameter $x$ is transcendental under a counterexample, so the usual classifications of algebraic hypergeometric functions with rational parameters do not apply directly.

\section{What the explored routes do and do not prove}
\label{sec:failures}

It is useful to record the failure modes explicitly.

\subsection{Reciprocal gamma zeros}
The quotient $Q_b$ gives an exact entire coprimality problem and removes all integral solutions.  Its zero geometry alone is the original logarithmic equation.  Nevanlinna theory controls densities of zeros, but one additional common zero is an arithmetic event too sparse for density estimates.

\subsection{Gamma and Hurwitz multiplication}
These formulas produce large towers of normalized algebraic values on denominator grids generated by $2$ and $3$.  Their cocycle law never leaves the rank-two group.  A fifth-base theorem is needed.

\subsection{Polylogarithmic holonomy}
The integral case is rational/D-finite/$q$-regular, and the nonintegral case is not.  However, the two integer Cartier eigenvalues do not imply finite holonomic or regular rank.  Treating the eigenrelations as Mahler equations would be an error: they are extraction identities, not closed equations involving only $F_x(z),F_x(z^q),\ldots$.

\subsection{Mellin/Barnes resonance}
The double pole is exactly equivalent to the equality of the two logarithmic exponents.  The resulting $t^{-x}\log(1/t)$ term has an integer-coefficient generating function behind it, but such singular exponents need not be integral.  Natural-boundary theorems do not constrain the pole abscissa.

\subsection{$q$-logarithm irrationality}
The two normalized $q$-digamma values are irrational.  Irrationality at each base separately is compatible with every genuine integral solution.  A cross-base algebraic independence theorem, not two one-base irrationality theorems, is required.

\subsection{Numerical near-resonance}
High-precision searches can find integer pairs $(M,N)$ for which $\log M/\log2$ and $\log N/\log3$ are extremely close.  In the Mellin picture this means two nearby simple poles that mimic a double pole over a finite scale range.  Numerical asymptotics therefore cannot certify exact resonance without a rigorous lower bound for a bilinear expression in logarithms, which is precisely the hard arithmetic core.

\section{Four settlement lemmas}
\label{sec:settlement}

The investigation isolates four sharply stated targets.  Any one would settle the conjecture.

\begin{criterion}[Special-function rank amplification]
Suppose that for real $x$ the algebraicity of $\mathfrak G_2(x)$ and $\mathfrak G_3(x)$ implies the algebraicity of $\mathfrak G_5(x)$.  Then $x\in\Z$ whenever $2^x,3^x\in\Z$.
\end{criterion}

\begin{proof}
The transfer quotients equal $2^x,3^x,5^x$.  If $x$ were nonintegral, \Cref{thm:spectrum} would force $5^x$ to be transcendental.
\end{proof}

\begin{criterion}[Cartier finite-rank lemma]
Suppose that a polylogarithm $F_x=\Li_{-x}$ which is a simultaneous $U_2,U_3$ eigenfunction with algebraic-integer eigenvalues is necessarily D-finite (or merely $q$-regular for one $q\ge2$).  Then the conjecture holds.
\end{criterion}

\begin{proof}
Apply \Cref{thm:holonomic} or \Cref{prop:qregular}.
\end{proof}

\begin{criterion}[Barnes resonance integrality]
Suppose that for positive integers $M,N$, a double real pole of \eqref{eq:mellin-transform} can occur only at an integer abscissa.  Then the conjecture holds.
\end{criterion}

\begin{proof}
By \Cref{thm:resonance}, the pole is at the common exponent $x$.
\end{proof}

\begin{criterion}[Double-sine quotient rigidity]
Let $\omega_1=\log2$, $\omega_2=\log3$, and $z_x$ be as in \eqref{eq:zx}.  Suppose algebraicity of both shift quotients \eqref{eq:S2shift1}--\eqref{eq:S2shift2} at $z_x$ forces $x\in\Q$.  Then the conjecture holds.
\end{criterion}

The first criterion is probably the most direct: it asks special-function arithmetic to manufacture a third independent exponential.  The second connects to fast-moving multivariate Mahler and Cobham theory.  The third is an analytic-combinatorial rigidity statement about an explicitly given integer series.  The fourth packages both bases into one two-period function and may be approachable through modular-double methods.

\section{A focused research program}
\label{sec:program}

The following sequence of tasks is designed to avoid rank-preserving loops.

\paragraph{1. Gamma/Hurwitz determinant experiments.}
Construct determinants from the denominator-$2^a3^b$ gamma or Hurwitz grids which vanish identically at integral $x$ because of finite polynomial or finite-difference rank.  Prove that a nonzero determinant would be algebraic under \eqref{eq:MN}.  If its analytic evaluation contains a $5^x$ factor, \Cref{crit:rank-amplification} finishes the argument.  The crucial requirement is that the determinant use an additive recombination of grid points, not only multiplication formulas.

\paragraph{2. Cartier modules with arithmetic coefficient fields.}
Study the $\Qbar$-span of congruence subsequences of $(n^x)$.  \Cref{cor:integer-spectrum} gives an unusually sharp coefficient-field dichotomy: all $3$-smooth coefficients are integers and all others are transcendental.  A theorem showing that simultaneous algebraic Cartier eigenvalues and this dichotomy force a finite module would imply integrality by \Cref{prop:qregular}.

\paragraph{3. Mixed Mahler values.}
Apply multivariate Mahler methods to $H_{M,N}$ at algebraic $z$ with $0<|z|<1$.  The equation \eqref{eq:factored-mahler} is singularly simple, and the transformations $z\mapsto z^2,z^3$ are multiplicatively independent.  The immediate goal is not merely transcendence of $H_{M,N}(z)$---which is compatible with both cases---but an algebraic-independence statement whose exceptional locus detects the resonance $\alpha=\beta$ and then constrains its location.

\paragraph{4. Hyperbolic gamma modular duality.}
Translate \eqref{eq:S2M}--\eqref{eq:S2N} into the hyperbolic gamma/quantum dilogarithm normalization, where modular-dual functional equations are most symmetric.  Search for a connection formula whose coefficient is an algebraic function of the two contiguous quotients but also equals $c^x$ for $c\notin\Div\langle2,3\rangle$.

\paragraph{5. Quantitative entire corona estimates.}
Use the explicit quotient \eqref{eq:Qb} and uniform Stirling estimates to seek a lower bound
\[
 |Q_2(z)|+|Q_3(z)|\ge \exp[-\exp(C|z|)]
\]
outside controlled disks around the known zero lattices.  On the real axis the remaining issue becomes a simultaneous lower bound for $\dist(2^t,\N)$ and $\dist(3^t,\N)$.  Even a bound valid on a strategically chosen sequence of intervals could combine with the functional equations to exclude a common zero.

\section{Literature audit through 21 August 2026}
\label{sec:literature}

A dated web and arXiv search was carried out for the exact conjecture and for the adjacent theories used here: Four/Six Exponentials, multivariate Mahler functions and base-change theorems, Barnes double gamma and hyperbolic gamma, Hurwitz-zeta algebraic independence, and arithmetic of generalized $q$-logarithms.  No paper directly claiming an unconditional proof of the stated $2^x,3^x$ conjecture was located.  This is a search report, not a claim of exhaustive nonexistence.

The most relevant continuing developments are methodological rather than direct:
\begin{itemize}
\item multivariate Mahler theory and Cobham-type rationality results provide a framework for simultaneous transformations $z\mapsto z^2,z^3$ \cite{AdamczewskiBell,AdamczewskiFaverjonI,AdamczewskiFaverjonII};
\item modern treatments of Barnes multiple zeta/gamma and multiple sine functions give canonical two-period analytic objects and modular transformations \cite{Ruijsenaars2000,Narukawa2004};
\item Pad\'e and Hankel methods for Lambert series and $q$-logarithms give unconditional irrationality statements such as \Cref{thm:qdigamma-irr} \cite{Borwein1992,BundschuhVaananen1994};
\item the core transcendence obstruction remains a rank-two versus rank-three exponential problem, naturally adjacent to the Four Exponentials Conjecture and controlled unconditionally only after adding a third independent logarithm via the Six Exponentials Theorem \cite{Lang1966,Waldschmidt2000}.
\end{itemize}

Appendix~\ref{app:arxiv} records a machine-generated sample of recent arXiv metadata from the search queries.  It is included for reproducibility and should not be read as a claim that every returned preprint bears directly on the conjecture.

\section{Conclusion}
\label{sec:conclusion}

The special-function approach yields a coherent picture.

First, a hypothetical counterexample has an exact algebraic-base spectrum: the exponentiation map $a\mapsto a^x$ sends algebraic numbers back to algebraic numbers only on the divisible rank-two group generated by $2$ and $3$.  This is an unconditional consequence of Six Exponentials and is the strongest reusable theorem in the report.

Second, gamma, Hurwitz zeta, Lerch, and polylogarithmic multiplication formulas all expose the same character $q\mapsto q^x$.  They make the algebraic spectrum visible but do not enlarge it.  Their limitation is now precise: they are rank preserving.

Third, the mixed Mahler series \eqref{eq:HMN} turns the equality of the two logarithmic exponents into a genuine Barnes double-zeta resonance.  The double pole, logarithmic main term, and two incommensurable periodic fluctuations are exact.  The correct Barnes object lies in this Mellin kernel, not in a nonexistent two-dimensional common divisor of reciprocal gamma functions.

Fourth, $q$-gamma and $q$-digamma theory gives nontrivial arithmetic consequences at the unknown $x$, including two irrational generalized $q$-logarithms.  Hypergeometric addition gives the opposite phenomenon: values at bases $2$ and $3$ are integers while the corresponding value at base $5$ must be transcendental under a counterexample.

The missing step is therefore not another multiplication identity.  It is rank amplification: an argument which uses the simultaneous special-function structures at $2$ and $3$ to force one algebraic value outside $\Div\langle2,3\rangle$, or an independent finite-rank theorem that collapses $\Li_{-x}$ to integral order.  The four settlement lemmas in \Cref{sec:settlement} isolate concrete forms that such a breakthrough could take.

\appendix
\section{Supplementary proofs and checks}
\label{app:proofs}

\subsection{Cancellation in the entire quotient}
Let $k\in\Z_{\ge0}$.  Near $z=k$,
\[
 \Gamma(-z)\sim\frac{(-1)^{k+1}}{k!(z-k)}.
\]
Put $R=b^k-1$.  Since $1-b^z=-R-b^k\log b\,(z-k)+O((z-k)^2)$ and
\[
 \Gamma(w)\sim\frac{(-1)^R}{R!(w+R)}\qquad(w\to-R),
\]
one obtains
\[
 Q_b(k)=(-1)^{k+R}\frac{R!\,b^k\log b}{k!}\ne0.
\]
This verifies the removable singularity and nonvanishing asserted in \Cref{prop:Qb}.

\subsection{Linear independence of shifted powers}
Let $\lambda\notin\Z_{\ge0}$ and let $\alpha_1,\ldots,\alpha_r$ be distinct.  Suppose
\[
 \sum_{j=1}^r c_j(n+\alpha_j)^\lambda=0
\]
for all sufficiently large positive integers $n$.  Dividing by $n^\lambda$ and expanding in $1/n$ gives
\[
 \binom{\lambda}{m}\sum_{j=1}^r c_j\alpha_j^m=0\qquad(m\ge0).
\]
All generalized binomial coefficients are nonzero, and the equations for $m=0,\ldots,r-1$ are a Vandermonde system.  Hence every $c_j=0$.  This proves the independence used in both \Cref{thm:holonomic} and \Cref{prop:qregular}.

\subsection{Pole coefficients in the resonant Mellin transform}
With $v=s-x$, $\omega_1=\log2$, $\omega_2=\log3$,
\[
 \frac{\Gamma(x+v)}{(1-\e^{-\omega_1v})(1-\e^{-\omega_2v})}
 =\frac{A}{v^2}+\frac{B}{v}+O(1),
\]
where
\[
 \Gamma(x+v)=\Gamma(x)(1+\psi(x)v+O(v^2))
\]
and
\[
 \frac1{(1-\e^{-\omega_1v})(1-\e^{-\omega_2v})}
 =\frac1{\omega_1\omega_2v^2}
 +\left(\frac1{2\omega_1}+\frac1{2\omega_2}\right)\frac1v+O(1).
\]
Multiplication gives \eqref{eq:A}--\eqref{eq:B}.  Since
$t^{-x-v}=t^{-x}\exp(v\log(1/t))$, the residue of
$(A/v^2+B/v)t^{-x-v}$ is
$t^{-x}[A\log(1/t)+B]$.

\subsection{Why the oscillatory residue sums converge}
At the $k$th nonreal $2$-pole, Stirling's formula gives
\[
 |\Gamma(x+2\pi i k/\log2)|
 \ll_x (1+|k|)^{x-1/2}\exp\!\left(-\frac{\pi^2|k|}{\log2}\right).
\]
The other denominator is proportional to the distance of
$k\log3/\log2$ from an integer.  A lower bound for the nonzero linear form
$k\log3-\ell\log2$ is polynomial in the height of $(k,\ell)$ by Baker's theory.  Thus the reciprocal denominator grows at most polynomially, while the gamma factor decays exponentially.  This proves absolute convergence of \eqref{eq:P2}; the argument for \eqref{eq:P3} is identical.

\section{Dated arXiv search sample}
\label{app:arxiv}
\paragraph{Four exponentials and related algebraic independence.}
\begin{itemize}
\item Metadata query error recorded by the build script: \texttt{URLError(gaierror(-3, 'Temporary failure in name resolution'))}.
\end{itemize}
\paragraph{Multivariate and multiplicatively independent Mahler systems.}
\begin{itemize}
\item Metadata query error recorded by the build script: \texttt{URLError(gaierror(-3, 'Temporary failure in name resolution'))}.
\end{itemize}
\paragraph{Barnes multiple gamma and related functions.}
\begin{itemize}
\item Metadata query error recorded by the build script: \texttt{URLError(gaierror(-3, 'Temporary failure in name resolution'))}.
\end{itemize}
\paragraph{Arithmetic of q-logarithms.}
\begin{itemize}
\item Metadata query error recorded by the build script: \texttt{URLError(gaierror(-3, 'Temporary failure in name resolution'))}.
\end{itemize}
\paragraph{Hurwitz zeta arithmetic.}
\begin{itemize}
\item Metadata query error recorded by the build script: \texttt{URLError(gaierror(-3, 'Temporary failure in name resolution'))}.
\end{itemize}

\begin{thebibliography}{99}
\bibitem{AE44}
L. Alaoglu and P. Erd\H{o}s,
\emph{On highly composite and similar numbers},
Trans. Amer. Math. Soc. \textbf{56} (1944), 448--469.

\bibitem{Lang1966}
S. Lang,
\emph{Introduction to Transcendental Numbers},
Addison--Wesley, 1966.

\bibitem{Waldschmidt2000}
M. Waldschmidt,
\emph{Diophantine Approximation on Linear Algebraic Groups: Transcendence Properties of the Exponential Function in Several Variables},
Grundlehren Math. Wiss. 326, Springer, 2000.

\bibitem{WaldschmidtOpen}
M. Waldschmidt,
\emph{Open Diophantine problems},
Moscow Math. J. \textbf{4} (2004), 245--305;
\url{https://arxiv.org/abs/math/0312440}.

\bibitem{Baker1966}
A. Baker,
\emph{Linear forms in the logarithms of algebraic numbers I--IV},
Mathematika \textbf{13--15} (1966--1968).

\bibitem{Stanley1980}
R. P. Stanley,
\emph{Differentiably finite power series},
European J. Combin. \textbf{1} (1980), 175--188.

\bibitem{FlajoletGerholdSalvy}
P. Flajolet, S. Gerhold, and B. Salvy,
\emph{On the non-holonomic character of logarithms, powers, and the $n$th prime function},
Electron. J. Combin. \textbf{11}(2) (2005), Article A2;
\url{https://arxiv.org/abs/math/0501379}.

\bibitem{AlloucheShallit1992}
J.-P. Allouche and J. Shallit,
\emph{The ring of $k$-regular sequences},
Theoret. Comput. Sci. \textbf{98} (1992), 163--197.

\bibitem{AdamczewskiBell}
B. Adamczewski and J. P. Bell,
\emph{A problem about Mahler functions},
Ann. Sc. Norm. Super. Pisa Cl. Sci. (5) \textbf{17} (2017), 1301--1355;
\url{https://arxiv.org/abs/1303.2019}.

\bibitem{AdamczewskiFaverjonI}
B. Adamczewski and C. Faverjon,
\emph{Mahler's method in several variables I: The theory of regular singular systems},
preprint;
\url{https://arxiv.org/abs/1809.04823}.

\bibitem{AdamczewskiFaverjonII}
B. Adamczewski and C. Faverjon,
\emph{Mahler's method in several variables II: Applications to base change problems and finite automata},
preprint;
\url{https://arxiv.org/abs/1809.04826}.

\bibitem{Nishioka1996}
K. Nishioka,
\emph{Mahler Functions and Transcendence},
Lecture Notes in Math. 1631, Springer, 1996.

\bibitem{Polya1921}
G. P\'olya,
\emph{Arithmetische Eigenschaften der Reihenentwicklungen rationaler Funktionen},
J. reine angew. Math. \textbf{151} (1921), 1--31.

\bibitem{Lech1953}
C. Lech,
\emph{A note on recurring series},
Ark. Mat. \textbf{2} (1953), 417--421.

\bibitem{Barnes1901}
E. W. Barnes,
\emph{The theory of the double gamma function},
Philos. Trans. Roy. Soc. London Ser. A \textbf{196} (1901), 265--387.

\bibitem{Ruijsenaars2000}
S. N. M. Ruijsenaars,
\emph{On Barnes' multiple zeta and gamma functions},
Adv. Math. \textbf{156} (2000), 107--132.

\bibitem{Kurokawa1991}
N. Kurokawa,
\emph{Multiple sine functions and Selberg zeta functions},
Proc. Japan Acad. Ser. A Math. Sci. \textbf{67} (1991), 61--64.

\bibitem{Narukawa2004}
A. Narukawa,
\emph{The modular properties and the integral representations of the multiple elliptic gamma functions},
Adv. Math. \textbf{189} (2004), 247--267;
\url{https://arxiv.org/abs/math/0306164}.

\bibitem{Borwein1992}
P. B. Borwein,
\emph{On the irrationality of certain series},
Math. Proc. Cambridge Philos. Soc. \textbf{112} (1992), 141--146.

\bibitem{BundschuhVaananen1994}
P. Bundschuh and K. V\"a\"an\"anen,
\emph{Arithmetical investigations of a certain infinite product},
Compositio Math. \textbf{91} (1994), 175--199.

\bibitem{GasperRahman}
G. Gasper and M. Rahman,
\emph{Basic Hypergeometric Series}, 2nd ed.,
Cambridge University Press, 2004.

\bibitem{Lewin1981}
L. Lewin,
\emph{Polylogarithms and Associated Functions},
North-Holland, 1981.

\bibitem{Erdelyi1953}
A. Erd\'elyi et al.,
\emph{Higher Transcendental Functions}, Vol. I,
McGraw--Hill, 1953.
\end{thebibliography}

\end{document}
